\documentclass[../main]{subfiles}
\ifSubfilesClassLoaded{\setcounter{chapter}{7}}{}
\begin{document}

\chapter{Protokol dan Manajemen Kunci}

\begin{subcpmk}
  \item Sub-CPMK 6.1: Menganalisis protokol Diffie-Hellman dan siklus hidup manajemen kunci
\end{subcpmk}

% ============================================================
% MATERI POKOK
% ============================================================
\section{PKI, Chain of Trust, dan Protokol TLS}

\subsection{Peran Diffie-Hellman dalam TLS}

Diffie-Hellman memungkinkan dua pihak menyepakati kunci rahasia melalui kanal tidak aman \cite{stallings2017}. Parameter publik: prima $p$, generator $g$. Keamanan bergantung pada kesulitan logaritma diskrit (diperkenalkan di Bab 6).

\begin{enumerate}
  \item Alice memilih $a$ rahasia, mengirim $A = g^a \bmod p$ ke Bob
  \item Bob memilih $b$ rahasia, mengirim $B = g^b \bmod p$ ke Alice
  \item Alice menghitung $K = B^a \bmod p$
  \item Bob menghitung $K = A^b \bmod p$
  \item Keduanya memperoleh $K$ yang sama
\end{enumerate}

\subsection{TLS (Transport Layer Security)}

TLS mengamankan komunikasi di internet (HTTPS) \cite{stallings2017}. Menggabungkan pertukaran kunci (DH atau ECDH), autentikasi sertifikat, dan enkripsi simetris (AES-GCM, ChaCha20-Poly1305). Handshake: negosiasi cipher suite, verifikasi sertifikat server, pertukaran kunci sesi, lalu enkripsi aplikasi.

\begin{itemize}
  \item \textbf{Negosiasi cipher suite}: Server dan klien menegosiasi algoritma yang didukung.
  \item \textbf{Autentikasi server}: Server membuktikan identitas dengan sertifikat digital.
  \item \textbf{Pertukaran kunci}: Diffie-Hellman atau RSA untuk pertukaran kunci sesi.
  \item \textbf{Enkripsi data}: AES-GCM atau ChaCha20-Poly1305 untuk data.
  \item \textbf{Integritas}: AEAD (GCM) atau HMAC untuk integritas.
\end{itemize}

\subsection{Chain of Trust dan Sertifikat}

TLS membutuhkan \textbf{otentikasi server} melalui sertifikat X.509 \cite{stallings2017}. \textbf{Chain of trust}: sertifikat server ditandatangani CA menengah; CA menengah ditandatangani CA root. Klien memvalidasi rantai sampai CA root dalam daftar tepercaya (trust store). Tanpa validasi sertifikat, koneksi rentan man-in-the-middle.

TLS 1.2 dan TLS 1.3 adalah versi saat ini. Forward secrecy: kunci sesi ephemeral sehingga kebocoran kunci privat tidak membeberkan sesi lampau \cite{stallings2017}.

\section{Manajemen Kunci}

\subsection{Siklus Hidup Kunci}

Manajemen kunci kriptografi modern memerlukan pendekatan yang komprehensif untuk memastikan keamanan, ketersediaan, dan kepatutan terhadap regulasi. Menurut Microsoft, desain PKI yang baik harus mempertimbangkan berbagai aspek operasional, keamanan, dan kepatuhan bisnis. Manajemen kunci yang efektif memerlukan kombinasi teknologi, prosedur, dan kebijakan yang terintegrasi.

\begin{itemize}
  \item \textbf{Generasi}: Kunci harus dihasilkan dengan sumber acak yang aman (CSPRNG)
  \item \textbf{Distribusi}: Kunci simetris dikirim melalui kanal aman atau pertukaran Diffie-Hellman
  \item \textbf{Penyimpanan}: Kunci privat disimpan terenkripsi, dilindungi password atau HSM
  \item \textbf{Penggunaan}: Hindari penggunaan ulang kunci untuk tujuan berbeda
  \item \textbf{Pembaruan}: Rotasi kunci secara berkala
  \item \textbf{Pencabutan}: Jika kunci bocor, cabut dan ganti segera
\end{itemize}

Manajemen kunci memerlukan pendekatan yang sistematis dan prosedural yang ketat. Menurut sumber keamanan, setiap aspek manajemen harus didokumentasikan dan diaudit secara berkala untuk memastikan konsistensi dan kepatuhan. Implementasi manajemen kunci yang buruk dapat menyebabkan kebocoran kunci, downtime sistem, atau bahkan kegagalan keamanan perusahaan.

\subsection{Distribusi Kunci}

Untuk kriptografi simetris, distribusi kunci merupakan tantangan yang kompleks. Solusi modern menggunakan arsitektur PKI (Public Key Infrastructure) yang terstandar. Menurut Microsoft, PKI menyediakan kerangka kerja yang terpisah untuk Certificate Authority (CA), Registration Authority (RA), dan Validation Authority (VA), memungkinkan manajemen siklus hidup kunci yang terpusat dan skalabel.

PKI terdiri dari beberapa komponen fundamental yang bekerja sama untuk memastikan keamanan:
\begin{itemize}
  \item \textbf{Certificate Authority (CA)}: Menerbitkan dan menandatangani sertifikat digital
  \item \textbf{Registration Authority (RA)}: Mendaftarkan identitas pemohon sertifikat
  \item \textbf{Validation Authority (VA)}: Memvalidasi permohonan sertifikat
  \item \textbf{Certificate Database}: Menyimpan dan mengelola sertifikat yang diterbitkan
  \item \textbf{Certificate Revocation}: Mencabut dan mendistribusikan daftar sertifikat yang dicabut
  \item \textbf{Key Management}: Menyimpan dan melindungi kunci CA
\end{itemize}

Arsitektur PKI yang dirancang dengan baik memungkinkan segregasi tugas, delegasi otoritas, dan redundansi. Menurut sumber keamanan, implementasi PKI yang sukses harus mengikuti standar industri seperti X.509 untuk format sertifikat dan protokol manajemen kunci yang aman seperti KMIP (Key Management Interoperability Protocol).

Untuk kriptografi simetris, distribusi kunci merupakan tantangan tersendiri. Solusi praktis meliputi: pertukaran Diffie-Hellman untuk menghasilkan kunci sesi; enkripsi kunci simetris dengan kunci publik penerima (hybrid encryption); serta Key Distribution Center (KDC) dalam skema terpusat. Pertukaran Diffie-Hellman tetap menjadi pilihan utama dalam protokol modern seperti TLS karena efisiensi komputasional dan forward secrecy. Enkripsi kunci simetris dengan kunci publik penerima menyediakan keamanan end-to-end yang baik. KDC memungkinkan kontrol terpusat atas kunci dan kebijakan yang konsisten untuk organisasi besar.


% ============================================================
% AKTIVITAS PEMBELAJARAN
% ============================================================

\begin{aktivitas}
  \item \textbf{Analisis}: Jalankan protokol Diffie-Hellman secara manual dengan $p=23$, $g=5$, $a=6$, $b=15$. Verifikasi bahwa Alice dan Bob memperoleh $K$ yang sama.
  
  \item \textbf{Studi Kasus}: Analisis bagaimana HTTPS (TLS) melindungi transaksi e-commerce.
  
  \item \textbf{Diskusi}: Apa risiko jika kunci privat tidak disimpan dengan aman?
\end{aktivitas}

% ============================================================
% LATIHAN DAN REFLEKSI
% ============================================================

\begin{latihan}
  \item Jelaskan langkah protokol Diffie-Hellman. Mengapa Eve tidak dapat memperoleh $K$?
  
  \item Sebutkan fase-fase dalam siklus hidup kunci.
  
  \item Bagaimana TLS menggabungkan kriptografi simetris dan asimetris?
  
  \item Apa peran sertifikat digital dalam TLS?
\end{latihan}

% ============================================================
% ASESMEN
% ============================================================

\begin{asesmen}
Jelaskan dalam essay: (a) Mengapa Diffie-Hellman aman terhadap eavesdropping? (b) Apa yang terjadi jika kunci tidak dirotasi?
\end{asesmen}

% ============================================================
% CHECKLIST KOMPETENSI
% ============================================================

\begin{checklist}
  \item Saya dapat menjelaskan protokol Diffie-Hellman
  \item Saya dapat menganalisis siklus hidup manajemen kunci
  \item Saya memahami peran TLS dalam keamanan komunikasi
  \item Saya dapat menjelaskan tantangan distribusi kunci
\end{checklist}

% ============================================================
% RANGKUMAN
% ============================================================

\begin{rangkuman}
Bab ini membahas protokol Diffie-Hellman untuk pertukaran kunci, TLS untuk keamanan transport, serta manajemen kunci (siklus hidup, distribusi).

\textbf{Poin Kunci:}
\begin{itemize}
  \item Diffie-Hellman: pertukaran kunci aman tanpa pertukaran rahasia
  \item TLS: negosiasi, autentikasi, pertukaran kunci, enkripsi data
  \item Manajemen kunci: generasi, distribusi, penyimpanan, rotasi, pencabutan
\end{itemize}

\textbf{Kata Kunci}: \crypt{Diffie-Hellman}, \crypt{TLS}, \crypt{Manajemen Kunci}, \crypt{Sertifikat}
\end{rangkuman}

\ifSubfilesClassLoaded{
  \renewcommand{\bibname}{Daftar Pustaka}
  \bibliographystyle{plain}
  \bibliography{../references}
}{}
\end{document}
